% Paper 10: The Crypto Fear Channel — Asymmetric BTC-Equity Volatility Spillover
% main.tex: canonical full body draft (renamed from body_v5.tex 2026-04-28;
%           v2.2 batch fix 2026-04-28 — 11 issues from v1 review closed)
% 9 sections, 16+ pages, 6 tables, 22 bibitems, ~4,900+ body words.
% Numbers all sourced from experiments/k1025/k1025_results.json
% (verified by reproduce.py 29/29 byte-match, alert_level=green).

\documentclass[11pt,a4paper]{article}
\usepackage{graphicx}
\usepackage{amsmath,amssymb}
\usepackage{natbib}
\usepackage[margin=1in]{geometry}
\usepackage{booktabs}
\usepackage{url}
\usepackage{float}

\title{The Crypto Fear Channel:\\ Asymmetric, Tail-Concentrated, and Regime-Dependent Volatility Spillover from Bitcoin to Equity Markets}

\author{Yi-Hao Lai\thanks{Department of Finance, Da-Yeh University, Taiwan. Email: yihao.lai@gmail.com.}}

\date{\today}

\begin{document}
\maketitle

\begin{abstract}
\noindent Using daily returns on SPY, BTC-USD, and VIX from 2015-02 to 2026-04 ($N = 2{,}812$), we document three stylized facts about the crypto-to-equity volatility spillover. First, the spillover is \emph{asymmetric}: Bitcoin downside volatility Granger-causes VIX at lags 1--5, whereas Bitcoin upside volatility does not (symmetric Granger tests covering lags 1--10 corroborate the broader directional result).% source: experiments/k1025/k1025_results.json .asymmetric_granger (lags 1-5) + .granger_causality.btc_rv_to_vix (symmetric 1-10)
 Second, the spillover is \emph{tail-concentrated with sign reversal}: the quantile-regression coefficient of VIX on BTC realized variance is significantly \emph{negative} at lower quantiles ($-2.86$ at $\tau = 0.05$; $-2.34$ at $\tau = 0.25$), indicating bull-market decoupling, then turns positive at the median ($+2.61$ at $\tau = 0.5$), continues climbing through the upper interquartile range ($+8.76$ at $\tau = 0.75$), and amplifies to $+22.31$ at the 95th percentile, yielding an $8.5\times$ upper-tail amplification relative to the median. Third, the spillover is \emph{regime-dependent}: in a five-subperiod breakdown, Granger causality is statistically significant only during 2020 ($F = 11.05$, $p < 0.001$) and non-significant in the other four subperiods (2015--2017, 2018--2019, 2021--2022, and 2023--2026). In the Diebold-Yilmaz framework, Bitcoin is a \emph{net receiver}, rehabilitating the narrative from ``crypto as a fear originator'' to ``crypto as a fear amplifier'' conditional on pre-existing equity stress. Out-of-sample forecasts, however, show that adding BTC realized volatility does not statistically improve an AR model of VIX (Diebold-Mariano $t = -0.98$, $p = 0.33$), highlighting that Granger causality is a necessary but not sufficient condition for forecastability. Taken together, our findings reframe the crypto-equity spillover as a \emph{crisis-time amplifier}: informative for ex-post risk attribution and margin design, but inadequate as an ex-ante predictor of equity fear. \\

\noindent\textbf{Keywords:} volatility spillover, Bitcoin, VIX, asymmetric causality, quantile regression, Diebold-Yilmaz, crypto-equity linkage \\
\textbf{JEL:} G11, G15, G17, C22
\end{abstract}

\section{Introduction}

Since the institutionalization of Bitcoin trading and the 2024 approval of U.S.\ spot BTC ETFs, the question of whether cryptocurrency markets transmit fear to traditional equity markets has acquired first-order practical importance for portfolio risk management, prudential supervision, and margin design. Yet the empirical literature has yielded an uneasy consensus. On one hand, correlation-based studies report rising BTC-equity co-movement in stress periods, suggesting that cryptocurrency is no longer a standalone asset class.\footnote{See, e.g., \citet{bouri2020}, \citet{corbet2018}, and \citet{matkovskyy2019} for the early post-2017 literature.} On the other hand, return-predictability tests often fail to find that Bitcoin-based signals improve forecasts of equity volatility at practical horizons. This tension motivates a central question: is the crypto-equity fear channel \emph{real but nonlinear}, \emph{real but regime-dependent}, or an artifact of tail-period correlation without predictive content?

We resolve this tension by decomposing the spillover along three dimensions that the prior literature has typically examined in isolation: directional asymmetry, tail dependence, and regime stability. Each of these dimensions yields a stylized fact that materially changes how one should interpret correlation-based evidence.

\paragraph{Asymmetry.} Using the asymmetric Granger-causality framework of \citet{hatemi2012}, we separate Bitcoin positive and negative realized volatility and test whether each Granger-causes VIX at lags 1 through 5.% source: experiments/k1025/k1025_results.json .asymmetric_granger.btc_neg_to_vix (lags 1-5)
 We find that the negative branch carries the entire spillover: Bitcoin downside volatility Granger-causes VIX at every lag tested (1--5), while the upside branch is uniformly non-significant across the same lag window. This is not merely a statistical artifact of a leverage effect --- the asymmetry is corroborated by standard symmetric Granger tests over the broader 1 through 10 lag range, where BTC-to-VIX causality is significant at lags 2--10.% source: experiments/k1025/k1025_results.json .granger_causality.btc_rv_to_vix (symmetric lags 1-10)
 Our interpretation is behavioral: retail-heavy crypto markets transmit \emph{fear}, not \emph{euphoria}, into the institutional equity complex.

\paragraph{Tail concentration with sign reversal.} Quantile regressions of VIX on BTC realized variance reveal a richer structure than a simple tail amplification: the slope \emph{reverses sign} across the VIX distribution. At the lower quantiles the coefficient is significantly \emph{negative} ($-2.86$ at $\tau = 0.05$; $-2.34$ at $\tau = 0.25$), consistent with a bull-market coexistence regime in which rising Bitcoin volatility coincides with \emph{compressing} equity fear. The slope turns positive at the median ($+2.61$ at $\tau = 0.5$) and amplifies sharply in the upper tail ($+22.31$ at $\tau = 0.95$), an approximately $8.5\times$ amplification of the upper tail relative to the median ($\tau = 0.95 / \tau = 0.5 \approx 8.54$).% source: experiments/k1025/k1025_results.json .quantile_regression (tau=0.05,0.25,0.5,0.95)
 The spillover is therefore not a uniform channel, nor even a simple ``flat-then-spikes'' tail phenomenon, but a regime with a \emph{directional switch} at the median of VIX: BTC volatility dampens equity fear in calm markets and amplifies it in panic markets. Conventional OLS or DCC-GARCH estimates --- which collapse this sign-reversing heterogeneity into a single slope --- systematically mischaracterize both the bull-market decoupling and the crisis-time amplification of crypto volatility on equity fear.

\paragraph{Regime dependence.} A five-subperiod breakdown reveals that Granger significance is concentrated entirely in 2020 ($F = 11.05$, $p < 10^{-6}$). Pre-mania (2015--2017), crypto-winter (2018--2019), bull-bear (2021--2022), and recovery-plus-ETF (2023--2026) all fail to reject non-causality. Combined with the finding from \citet{diebold2012} spillover-index analysis that Bitcoin is a \emph{net receiver} rather than net sender, this suggests that the spillover is not a permanent feature of cross-market architecture but a \emph{conditional amplification mechanism} that activates only during pre-existing equity stress.

\paragraph{The forecastability gap.} A natural next question is whether these three stylized facts translate into out-of-sample predictive power. They do not. Comparing an AR($p$) specification of VIX against one augmented with BTC realized volatility over an OOS window of 1,826 days (2019--2026), the Diebold-Mariano test returns $t = -0.98$, $p = 0.33$ --- statistically indistinguishable from zero, and below the \citet{harvey2016} $|t| > 3$ threshold for newly proposed factors. This result echoes a broader methodological point: \emph{Granger causality is a necessary but not sufficient condition for forecastability}. An informative signal that appears only in tail events and only during rare regimes may be too sparse to outperform a well-specified autoregressive benchmark on average over a long OOS sample.

We make three contributions to the literature. First, we document a new empirical pattern in the BTC-to-VIX channel: the conditional response of equity fear to Bitcoin realized variance \emph{reverses sign} between the lower and upper halves of the VIX distribution and amplifies by approximately $8.5\times$ from the median to the 95th percentile. The sign reversal --- statistically significant in both halves with $|t|$ exceeding 5 in every quantile tested --- is invisible to conditional-mean estimators and has not, to our knowledge, been documented in the cryptocurrency-equity spillover literature. Second, the COVID-2020 subperiod emerges as a structural watershed for the channel: of five non-overlapping subperiods spanning 2015 through 2026, only 2020 rejects non-causality at conventional levels, while a $-77$ percentage-point net Diebold-Yilmaz directional spillover identifies Bitcoin as a fear \emph{amplifier} (recipient), not originator (sender). The combination implies that the spillover is a \emph{crisis-time conditional amplifier} rather than a steady-state channel --- with direct margin-and-capital-requirement implications during regime transitions. Third, we adopt a joint reporting discipline that pairs each in-sample finding with an out-of-sample evaluation under the \citet{harvey2016} multiple-testing threshold; the resulting transparent OOS NULL ($t = -0.98$, well below $|t| > 3$) supplies a methodological template for the volatility-forecasting literature on when in-sample structure should and should not be expected to translate into point-forecast improvement \citep{conrad2020}. Each contribution rests on the four methodological building blocks --- asymmetric Granger causality, quantile regression for tail dependence, Diebold-Yilmaz directional spillover, and Diebold-Mariano forecast comparison --- developed in \S\ref{sec:methodology}; the novelty lies in their joint application to the BTC-equity fear channel and in the empirical patterns they collectively reveal.

The remainder of the paper is organized as follows. Section~\ref{sec:lit} surveys the crypto-equity spillover literature and positions our contribution. Section~\ref{sec:data} describes the data and preliminaries. Section~\ref{sec:methodology} details the four core methodological building blocks (asymmetric Granger, quantile regression, Diebold-Yilmaz, Diebold-Mariano), preceded by a symmetric-Granger baseline. Sections~\ref{sec:results} and~\ref{sec:robustness} present the main results and robustness checks. Section~\ref{sec:oos} reports the out-of-sample forecasting evaluation. Section~\ref{sec:discussion} discusses mechanisms and policy implications. Section~\ref{sec:conclusion} concludes.

\section{Literature Review}
\label{sec:lit}

We organize the related literature along three threads: (i) the empirical literature on crypto-equity volatility spillover and the safe-haven debate, (ii) the methodological literature underlying our four building blocks (asymmetric Granger, quantile regression, Diebold-Yilmaz, Diebold-Mariano), and (iii) the recent calls for honest out-of-sample evaluation in volatility forecasting. Each thread has progressed independently; our contribution is to combine them within a single framework that decomposes the spillover into asymmetry, tail concentration, regime conditionality, and predictive power.

\subsection{Crypto-equity volatility spillover and the safe-haven debate}

The empirical literature on Bitcoin's role in equity portfolios has evolved through three phases. The first phase, exemplified by \citet{corbet2018} and \citet{matkovskyy2019}, documented that cryptocurrency price movements respond meaningfully to macro-financial news and that contagion-style co-movement with equity markets emerges in stress periods, contradicting the early ``digital gold'' framing. The second phase, motivated by COVID-19 and reviewed in \citet{bouri2020}, \citet{conlon2020}, \citet{shahzad2019}, and \citet{klein2018}, sharpened the safe-haven question with mixed verdicts: wavelet- and DCC-based correlation analysis showed that Bitcoin's risk profile is closer to a high-volatility risky asset than to gold, and that its safe-haven properties are at best regime-dependent. The third phase, post-2020, has documented increasing institutional integration: \citet{yarovaya2022} report that information transmission between cryptocurrency and traditional asset markets intensified during the pandemic, and \citet{akyildirim2020} link implied-volatility shocks to cryptocurrency returns.

Three gaps motivate our analysis. First, most prior studies test \emph{symmetric} causality or \emph{symmetric} correlation, conflating the upside and downside volatility channels that we will show carry economically distinct information. Second, the tail dependence of the spillover is typically inferred indirectly from rolling DCC or wavelet coherence, rather than directly estimated through quantile regression. Third, sample windows often end at 2020 or include 2020 as a single observation, leaving the post-COVID regime-stability question unaddressed --- a gap that matters precisely because COVID-2020 will turn out to be the structural watershed of the spillover. Our analysis closes these gaps by replacing each symmetric or correlation-based summary with a direct, regime-aware, asymmetric counterpart, while extending the sample through April 2026.

\subsection{Methodological building blocks}

Our four building blocks each have a settled methodological literature. Asymmetric Granger causality, the technique behind our headline asymmetry result, was formalized by \citet{hatemi2012}, who decompose each series into positive and negative cumulative innovations and test their respective Granger predictability. Earlier nonparametric extensions of Granger causality \citep{diks2006} and asymmetric volatility-spillover tests \citep{hong2001} provide complementary diagnostic tools that we use as robustness checks. Each of these methods was developed to address the specific concern that linear Granger tests can mask sign-dependent transmission --- precisely the asymmetry the present paper documents in the BTC-to-VIX channel.

Quantile regression, our tool for tail dependence, originates with \citet{koenker1978} and has subsequently been extended into systemic-risk applications. Most relevant to our setting is \citet{adrian2016}'s CoVaR framework, in which the conditional value-at-risk of one institution given another's stress is estimated through quantile regression on returns. We adopt the quantile-regression framework rather than CoVaR \emph{per se} because our object of interest is the conditional response of equity \emph{volatility} (VIX), not the conditional value-at-risk of equity returns; but the underlying statistical engine is the same.

The Diebold-Yilmaz spillover index is now a canonical tool for measuring directional volatility transmission across asset classes. Its development spans \citet{diebold2009}, \citet{diebold2012}, and \citet{diebold2014network}, and a recent IMF policy note \citep{iyer2022} applies the framework specifically to cryptocurrency-equity spillovers. We follow the 2012 specification (variance decomposition of a generalized VAR) and report rolling 252-day spillover indices to assess directional connectedness over time.

Finally, our forecasting evaluation rests on the \citet{diebold1995} test of equal predictive accuracy and the \citet{harvey2016} multiple-testing threshold of $|t| > 3$ for newly proposed factors. The latter has become a methodological norm for newly proposed predictors in asset pricing, and we adopt it in the volatility-forecasting setting to avoid the inflated false-positive rate associated with conventional $|t| > 1.96$ inference under multiple testing.

\subsection{Honest out-of-sample evaluation in volatility forecasting}

A recurring tension in the volatility-forecasting literature is the gap between in-sample statistical significance and out-of-sample economic value. \citet{conrad2020}'s GARCH-MIDAS application demonstrates that incorporating macro-economic information through long-run components can deliver in-sample fit improvements that fail to translate into forecast accuracy gains. The cumulative evidence reviewed there motivates an evaluation discipline in which a new predictor must clear both (i) in-sample significance under the Harvey threshold and (ii) out-of-sample improvement under Diebold-Mariano, before being promoted to a structural claim. The present paper adopts this discipline and reports a transparent out-of-sample null --- BTC realized volatility fails to improve AR-VIX forecasts at $t = -0.98$ --- alongside its in-sample asymmetry, tail-amplification, and regime-stability findings. The combination is, to our knowledge, the first such honest joint reporting in the crypto-equity spillover literature.

\section{Data and Preliminaries}
\label{sec:data}

\subsection{Sample and variables}

Our sample consists of daily data on three series: the SPDR S\&P~500 ETF (SPY), Bitcoin in U.S.\ dollar terms (BTC-USD), and the CBOE Volatility Index (VIX). All series are obtained from Yahoo Finance via the \texttt{yfinance} package and pinned to a fixed snapshot CSV (\texttt{paper/crypto-fear-channel/data/spy\_btc\_usd\_vix\_2015-2026.csv}, 4{,}116 rows including BTC's seven-day calendar) generated on 2026-04-27 with \texttt{auto\_adjust=False}, retaining both raw close and adjusted close columns to insulate replication from retroactive dividend reconciliations. The full sample spans 2015-02-02 through 2026-04-08, yielding $N = 2{,}812$ trading-day observations after aligning BTC to the SPY/VIX trading calendar (we drop weekend BTC observations because VIX is undefined). The out-of-sample window used in \S\ref{sec:methodology}.5 begins 2019-01-01 and runs through 2026-04-08, providing 1{,}826 OOS observations --- spanning the COVID-19 shock (2020), the late-cycle bull market (2021), the 2022 bear, and the post-ETF recovery (2023--2026).% source: experiments/k1025/k1025_results.json .sample_period + .oos_period + .n_observations

\subsection{Constructed variables}

We work with three constructed series. Daily log returns are computed from adjusted close: $r^{\text{spy}}_t = \log(P^{\text{spy}}_t / P^{\text{spy}}_{t-1})$ and $r^{\text{btc}}_t = \log(P^{\text{btc}}_t / P^{\text{btc}}_{t-1})$. Realized volatility is constructed as a 22-day rolling standard deviation of log returns, annualized by $\sqrt{252}$: $\text{RV}^{(20)}_t = \text{std}_{[t-21,t]}(r) \cdot \sqrt{252}$, computed separately for SPY and BTC. The VIX is used in level (it is already an annualized measure, $\%$). For all forward-looking specifications we lag explanatory variables by one day (\texttt{signal.shift(1)} convention) to enforce causal predictability and rule out look-ahead.

For the asymmetric Granger tests of \S\ref{sec:methodology}.2 we follow \citet{hatemi2012}'s sign-decomposition approach with one adaptation: rather than test the cumulative-positive and cumulative-negative innovations of BTC realized volatility directly, we first decompose BTC \emph{returns} into positive and negative components ($r^{\text{btc},+}_t = \max(r^{\text{btc}}_t, 0)$ and $r^{\text{btc},-}_t = -\min(r^{\text{btc}}_t, 0)$) and then construct directional realized-volatility series from each component using the same 22-day rolling quadratic-mean as in the symmetric specification: $\text{RV}^{\text{btc},\pm,(20)}_t = \sqrt{\frac{1}{20}\sum_{j=t-19}^{t}(r^{\text{btc},\pm}_j)^2} \cdot \sqrt{252}$. We take first differences of these directional RV series to ensure stationarity in the Granger specification (the level series carry near-unit-root persistence in 20-day windows), and use the first-differenced series $\Delta \text{RV}^{\text{btc},\pm,(20)}_t$ as the explanatory variables in the asymmetric Granger tests. This adaptation preserves the key identification of \citet{hatemi2012} --- separate sign-conditional Granger tests --- while allowing the volatility-channel interpretation that motivates our application.

\subsection{Descriptive statistics and stationarity}

Table~\ref{tab:descriptive} reports descriptive statistics for the four primary series. BTC daily returns exhibit mean $0.23\%$, daily standard deviation $3.76\%$, and excess kurtosis $7.58$. SPY daily returns are an order of magnitude less volatile in standard deviation ($1.12\%$), but exhibit a higher excess kurtosis ($14.15$) in this sample window, which is concentrated in a handful of COVID-2020 and 2022 crisis days; the sample-period excess-kurtosis ranking does not reverse the well-documented unconditional-tail-thickness ranking of BTC over SPY but reflects the influence of a small set of extreme equity outliers within the 2015--2026 window. The VIX averages $18.4$ with maximum $82.7$ (March 2020) and minimum $9.1$, spanning the calm 2017--2019 vol regime, the COVID shock, the 2022 inflation regime, and the 2024--2026 normalization.

\begin{table}[H]
\centering
\caption{Descriptive statistics for daily series, 2015-02-02 to 2026-04-08 ($N = 2{,}812$).}
\label{tab:descriptive}
\begin{tabular}{lcccc}
\toprule
Series & Mean & Std.\ dev. & Skewness & Excess kurtosis \\
\midrule
$r^{\text{btc}}$ (\%) & 0.229 & 3.764 & $-$0.093 & 7.579 \\
$r^{\text{spy}}$ (\%) & 0.056 & 1.117 & $-$0.307 & 14.150 \\
VIX & 18.382 & 7.110 & --- & --- \\
$\text{RV}^{(20)}_{\text{btc}}$ & 0.542 & 0.257 & --- & --- \\
\bottomrule
\end{tabular}
\begin{flushleft}
\footnotesize VIX min/max are 9.14 / 82.69. $\text{RV}^{(20)}_{\text{btc}}$ min/max are 0.098 / 1.701 (annualized). Source: \texttt{experiments/k1025/k1025\_results.json .descriptive\_statistics}.
\end{flushleft}
\end{table}

Augmented Dickey-Fuller (ADF) stationarity tests reject the unit-root null for all three modeling series at conventional levels: $\text{RV}^{(20)}_{\text{btc}}$ ($t = -4.90$, $p < 0.001$, lags 20), VIX ($t = -5.71$, $p < 10^{-6}$, lags 9), and $\text{RV}^{(20)}_{\text{spy}}$ ($t = -4.46$, $p < 0.001$, lags 20).% source: experiments/k1025/k1025_results.json .adf_tests
 Stationarity is a precondition for the Granger and Diebold-Yilmaz analyses below, both of which assume covariance stationarity. We do not differentiate further; the analyses use the level series.

\section{Methodology}
\label{sec:methodology}

We employ four methodological building blocks. The first three (\S\ref{sec:methodology}.1--\ref{sec:methodology}.4) characterize \emph{in-sample} cross-asset structure along three dimensions: directional asymmetry, tail dependence, and regime stability. The fourth (\S\ref{sec:methodology}.5) tests \emph{out-of-sample} predictive content.

\subsection{Symmetric Granger causality (baseline)}
\label{subsec:granger_sym}

As a baseline we test whether BTC realized volatility Granger-causes VIX in a standard linear specification with lags $\ell = 1, \ldots, 10$. For each lag $\ell$ we estimate the unrestricted regression
\begin{equation}
\text{VIX}_t = \alpha_0 + \sum_{j=1}^{\ell} \alpha_j \text{VIX}_{t-j} + \sum_{j=1}^{\ell} \beta_j \text{RV}^{(20)}_{\text{btc}, t-j} + \varepsilon_t
\label{eq:granger_unrestricted}
\end{equation}
and the restricted regression that drops the $\beta_j$ terms. The joint null $H_0: \beta_1 = \cdots = \beta_\ell = 0$ is tested with a standard $F$-test under heteroskedasticity- and autocorrelation-consistent (HAC) standard errors. We use the Newey-West kernel with the automatic bandwidth selection rule of \citet{andrews1991}, which is the default specification of the \texttt{statsmodels.tsa.stattools.grangercausalitytests} routine that generates our $F$-statistics. We tabulate $F$-statistics and $p$-values across lags 1 through 10 and treat the symmetric Granger result as the linear-causality benchmark against which the asymmetric and tail-conditional specifications below are compared.

\subsection{Asymmetric Granger causality (Hatemi-J)}
\label{subsec:granger_asym}

Following the sign-decomposition approach of \citet{hatemi2012}, we test whether the directional RV series defined in \S\ref{sec:data}.2 (the first-differenced rolling-RV series constructed separately from positive-only and negative-only BTC daily returns) carry distinct Granger predictive content for VIX. The asymmetric Granger specifications are estimated as two separate regressions, one per branch:
\begin{equation}
\begin{aligned}
\text{VIX}_t &= \alpha^+_0 + \sum_{j=1}^{\ell} \alpha^+_j \text{VIX}_{t-j} + \sum_{j=1}^{\ell} \beta^{+}_j \, \Delta \text{RV}^{\text{btc},+,(20)}_{t-j} + \varepsilon^+_t, \\
\text{VIX}_t &= \alpha^-_0 + \sum_{j=1}^{\ell} \alpha^-_j \text{VIX}_{t-j} + \sum_{j=1}^{\ell} \beta^{-}_j \, \Delta \text{RV}^{\text{btc},-,(20)}_{t-j} + \varepsilon^-_t.
\end{aligned}
\label{eq:granger_asym}
\end{equation}
We test the two null hypotheses, $H_0^+: \beta^{+}_1 = \cdots = \beta^{+}_\ell = 0$ (BTC upside RV innovations do not Granger-cause VIX) and $H_0^-: \beta^{-}_1 = \cdots = \beta^{-}_\ell = 0$ (BTC downside RV innovations do not Granger-cause VIX), each with a standard $F$-test against the corresponding restricted regression. The Hatemi-J framework is appropriate when the suspected transmission channel is sign-dependent --- our prior, motivated by retail-driven crypto trading and the asymmetric leverage effect in equity volatility, is precisely such asymmetry. We test lags $\ell \in \{1, 2, 3, 4, 5\}$.\footnote{The lag window 1--5 is chosen to match the timescale at which retail-investor sentiment plausibly transmits across asset markets; the symmetric Granger test in \S\ref{subsec:granger_sym} extends to $\ell = 10$ to confirm robustness across longer horizons.}

\subsection{Quantile regression for tail dependence}
\label{subsec:qr}

Linear regression measures the conditional mean response of VIX to BTC RV; conditional on the documented heavy tails of both series and the sign-reversing structure we will document, the conditional mean is uninformative about regime-specific transmission. We follow \citet{koenker1978} and estimate quantile regressions of VIX on BTC realized variance at five representative quantiles $\tau \in \{0.05, 0.25, 0.50, 0.75, 0.95\}$:
\begin{equation}
Q_{\tau}(\text{VIX}_t \mid \text{RV}^{(20)}_{\text{btc}, t-1}) = \alpha_\tau + \beta_\tau \cdot \text{RV}^{(20)}_{\text{btc}, t-1},
\label{eq:qr}
\end{equation}
estimated via the standard Koenker-Bassett check-loss minimization. Standard errors and confidence intervals are obtained via 1{,}000-replication bootstrap of the regression sample. The tail-dependence ratio reported in our headline result is the upper-tail amplification $\beta_{0.95} / \beta_{0.50}$, which exceeds 8 in our sample. We also report intermediate quantiles in robustness checks. The conceptual link to \citet{adrian2016}'s CoVaR framework is direct: both rely on quantile regression to recover state-dependent cross-asset transmission, with the present paper applying the technique to volatility (VIX) rather than tail risk.

\subsection{Diebold-Yilmaz spillover index}
\label{subsec:dy}

To measure directional volatility transmission, we follow the \citet{diebold2012} variance-decomposition specification. We fit a generalized vector autoregression of order $p = 4$ to the 3-variable system $\{\text{VIX}, \text{RV}^{(20)}_{\text{spy}}, \text{RV}^{(20)}_{\text{btc}}\}$ and compute the $H$-step-ahead generalized forecast-error variance decomposition (FEVD) at $H = 10$ days. The directional spillover from variable $i$ to variable $j$ is the share of variable $j$'s $H$-step forecast-error variance attributable to shocks originating in variable $i$. The total spillover index, the ``to'' index, the ``from'' index, and the net spillover (to minus from) are computed in the standard generalized framework. To track time variation, we compute the same indices on a rolling 252-trading-day window, allowing direct visualization of how the BTC--VIX--SPY spillover structure evolves across the 2015--2026 sample.

\subsection{Out-of-sample forecasting framework}
\label{subsec:dm}

We test whether the in-sample asymmetry, tail-amplification, and regime-stability findings translate into out-of-sample predictive power for VIX. The base forecast is an autoregressive AR($p$) model with lag length selected by AIC on the in-sample window (2015-02 to 2018-12); the augmented forecast adds lagged BTC realized volatility:
\begin{equation}
\text{VIX}_t = \alpha_0 + \sum_{j=1}^{p} \alpha_j \text{VIX}_{t-j} + \gamma \cdot \text{RV}^{(20)}_{\text{btc}, t-1} + u_t.
\label{eq:oos_aug}
\end{equation}
Both models are re-estimated on a rolling basis, and one-day-ahead point forecasts are generated for each day in the OOS window 2019-01-01 to 2026-04-08 (1{,}826 observations). Forecast errors are evaluated under squared-error loss, and the difference in loss between the two specifications is tested with the \citet{diebold1995} test of equal predictive accuracy with the small-sample adjustment of \citet{harvey1997}, evaluated under the \citet{harvey2016} threshold $|t| > 3$ for newly proposed predictors --- a tighter discipline than the conventional $|t| > 1.96$ benchmark, motivated by the multiple-testing concerns reviewed in \citet{conrad2020}. Failure to reject the null at this threshold is interpreted as a transparent signal that, despite in-sample causal structure, BTC realized volatility does not improve VIX point forecasts on average over the long OOS sample.

\section{Main Results}
\label{sec:results}

This section presents the three core empirical findings: directional asymmetry of the BTC-to-VIX channel (\S\ref{subsec:r_asym}), tail concentration with sign reversal in the conditional response (\S\ref{subsec:r_qr}), and regime-dependence of the Granger structure (\S\ref{subsec:r_regime}).

\subsection{Asymmetric Granger causality: only BTC downside transmits}
\label{subsec:r_asym}

Table~\ref{tab:granger_asym} reports the central asymmetric-Granger result. The negative branch ($\text{RV}^{\text{btc},-}$, the cumulative downside innovation) Granger-causes VIX at every tested lag from 1 through 5, with $F$-statistics ranging from 6.64 to 18.96 and $p$-values uniformly below $10^{-5}$. The positive branch ($\text{RV}^{\text{btc},+}$) is not significant at any lag, with $F$-statistics between 0.17 and 2.00 and $p$-values between 0.16 and 0.95. The asymmetry is monolithic: there is no overlap region in which both branches contribute --- the channel is fully one-sided.% source: experiments/k1025/k1025_results.json .asymmetric_granger.btc_pos_to_vix + .btc_neg_to_vix (lags 1-5)

\begin{table}[H]
\centering
\caption{Asymmetric Granger causality: BTC realized-volatility innovations to VIX, lags 1--5.}
\label{tab:granger_asym}
\begin{tabular}{lcccc}
\toprule
 & \multicolumn{2}{c}{BTC$^+$ $\to$ VIX} & \multicolumn{2}{c}{BTC$^-$ $\to$ VIX} \\
\cmidrule(lr){2-3} \cmidrule(lr){4-5}
Lag & $F$ & $p$ & $F$ & $p$ \\
\midrule
1 & 2.00 & 0.157 & 18.96 & $1.4 \times 10^{-5}$ \\
2 & 0.20 & 0.816 & 14.79 & $4.1 \times 10^{-7}$ \\
3 & 0.20 & 0.894 & 10.18 & $1.2 \times 10^{-6}$ \\
4 & 0.17 & 0.954 & 7.34  & $7.1 \times 10^{-6}$ \\
5 & 0.28 & 0.927 & 6.64  & $3.7 \times 10^{-6}$ \\
\bottomrule
\end{tabular}
\begin{flushleft}
\footnotesize $F$-tests of the joint null that the lagged BTC$^{\pm}$ coefficients are zero in Eq.~(\ref{eq:granger_asym}). The negative-branch test rejects the null at every lag with $p < 10^{-5}$; the positive branch fails to reject at any lag. Source: \texttt{experiments/k1025/k1025\_results.json .asymmetric\_granger}.
\end{flushleft}
\end{table}

The asymmetry is corroborated by symmetric Granger tests over a broader lag window. Pooling the positive and negative innovations into a single specification, the symmetric BTC RV $\to$ VIX test rejects the joint null at lags 2 through 10 with $p$-values from $0.021$ down to $10^{-7}$; the lag-1 specification is non-significant ($F = 0.64$, $p = 0.42$).% source: experiments/k1025/k1025_results.json .granger_causality.btc_rv_to_vix (lags 1-10)


\paragraph{Asymmetric vs.\ symmetric Granger: reading the lag-1 result.} The lag-1 contrast is the cleanest diagnostic in our test battery for why asymmetric decomposition matters. The symmetric specification at lag 1 averages over the (small, possibly zero) upside coefficient and the (large, positive) downside coefficient, producing the pooled $F = 0.64$ ($p = 0.42$) non-rejection. Once the two branches are estimated separately as in Eq.~(\ref{eq:granger_asym}), the downside-branch test recovers $F = 18.96$ ($p < 10^{-5}$): the lag-1 BTC-to-VIX channel is real but mechanically invisible to a sign-pooled symmetric test. Reading the two tables jointly: BTC RV does Granger-cause VIX, but the entire signal lives in the downside half of the BTC innovation distribution. This is a clean illustration of why asymmetric and symmetric Granger are not interchangeable diagnostics; using only the symmetric test would have led us to mis-time the predictive horizon by one lag.

\subsection{Tail dependence with sign reversal}
\label{subsec:r_qr}

Table~\ref{tab:qr} reports the quantile-regression slope of VIX on lagged BTC realized variance across five representative quantiles. The most striking pattern is the \emph{sign reversal} of the coefficient as the conditional VIX quantile rises: at $\tau = 0.05$ and $\tau = 0.25$ the slope is significantly \emph{negative} ($\beta_{0.05} = -2.86$, $t = -10.93$; $\beta_{0.25} = -2.34$, $t = -6.54$), indicating that in the lower (calm) tail of the VIX distribution, higher BTC realized variance coincides with \emph{lower} expected VIX. The slope crosses zero between $\tau = 0.25$ and $\tau = 0.5$ and turns sharply positive: $\beta_{0.5} = 2.61$ ($t = 5.44$), $\beta_{0.75} = 8.76$ ($t = 11.17$), $\beta_{0.95} = 22.31$ ($t = 8.88$).% source: experiments/k1025/k1025_results.json .quantile_regression (tau=0.05,0.25,0.5,0.75,0.95)

\begin{table}[H]
\centering
\caption{Quantile regression of VIX on lagged BTC realized variance, $N = 2{,}812$.}
\label{tab:qr}
\begin{tabular}{lccccc}
\toprule
$\tau$ & 0.05 & 0.25 & 0.50 & 0.75 & 0.95 \\
\midrule
$\beta_\tau$ & $-2.86$ & $-2.34$ & $+2.61$ & $+8.76$ & $+22.31$ \\
SE          & 0.262   & 0.358   & 0.480   & 0.784   & 2.512   \\
$t$-stat    & $-10.93$ & $-6.54$ & $5.44$  & $11.17$ & $8.88$  \\
$p$-value   & $< 10^{-26}$ & $< 10^{-10}$ & $< 10^{-7}$ & $< 10^{-27}$ & $< 10^{-18}$ \\
\bottomrule
\end{tabular}
\begin{flushleft}
\footnotesize Slope coefficients from quantile regressions of VIX on lagged BTC realized variance. All coefficients are statistically significant at the $10^{-7}$ level or better. The sign of the slope reverses between the lower quartile and the median; the upper-tail amplification ratio is $\beta_{0.95} / \beta_{0.5} \approx 8.54$. Source: \texttt{experiments/k1025/k1025\_results.json .quantile\_regression}.
\end{flushleft}
\end{table}

The economically meaningful summary statistic is the upper-tail amplification ratio $\beta_{0.95} / \beta_{0.5} \approx 8.54$: the conditional sensitivity of VIX to BTC realized variance is roughly 8.5 times stronger at the 95th percentile of VIX than at the median. The sign reversal at lower quantiles is, to our knowledge, a new finding for crypto-equity spillover: it implies that conditional-mean estimators (OLS, DCC-GARCH) average over a regime in which BTC volatility \emph{dampens} expected equity fear (the bull-market decoupling we discuss in \S\ref{sec:discussion}) and a regime in which it \emph{amplifies} fear, producing a pooled slope that mischaracterizes both halves.

\subsection{Regime dependence: COVID-2020 as structural watershed}
\label{subsec:r_regime}

We split the 2015--2026 sample into five non-overlapping subperiods --- pre-mania (2015--2017), crypto-winter (2018--2019), COVID (2020), bull-bear (2021--2022), and recovery-plus-ETF (2023--2026) --- and re-estimate the BTC-to-VIX Granger test within each subperiod with AIC-selected lag length. Table~\ref{tab:subperiod} reports the result. Granger causality is statistically significant only during 2020 ($F = 11.05$, $p = 7.9 \times 10^{-7}$, lag 3); each of the other four subperiods fails to reject the non-causality null at any conventional level.% source: experiments/k1025/k1025_results.json .subperiod_granger

\begin{table}[H]
\centering
\caption{Subperiod Granger tests of BTC realized volatility on VIX, AIC-best lag.}
\label{tab:subperiod}
\begin{tabular}{lcccc}
\toprule
Subperiod & $N$ & Best lag & $F$ & $p$ \\
\midrule
2015--2017 (Pre-mania)         & 735 & 1 & 0.59  & 0.443 \\
2018--2019 (Crypto winter)     & 503 & 1 & 0.23  & 0.630 \\
2020 (COVID)                   & 253 & 3 & 11.05 & $7.9 \times 10^{-7}$ \\
2021--2022 (Bull-Bear)         & 503 & 1 & 1.95  & 0.163 \\
2023--2026 (Recovery+ETF)      & 818 & 3 & 0.46  & 0.709 \\
\bottomrule
\end{tabular}
\begin{flushleft}
\footnotesize Granger $F$-tests within each subperiod with AIC-selected lag length. Only 2020 rejects the non-causality null. Source: \texttt{experiments/k1025/k1025\_results.json .subperiod\_granger}.
\end{flushleft}
\end{table}

The DCC correlation by VIX regime corroborates the regime-dependence finding from a different angle. Sorting trading days into four VIX regimes (Low: VIX $\le$ 15; Normal: 15 $<$ VIX $\le$ 22; High: 22 $<$ VIX $\le$ 30; Crisis: VIX $>$ 30), the mean BTC-SPY DCC correlation rises near-monotonically from 0.07 in the Low regime to 0.27 in Normal, peaking at 0.45 in High, and remaining elevated at 0.41 in Crisis (Table~\ref{tab:dcc}; the slight Crisis dip from 0.45 to 0.41 reflects the small Crisis-regime sample of $n = 63$). The percentage of days with positive correlation rises from 62\% in Low to 95\% in High and 98\% in Crisis. Bitcoin's purported diversification value relative to equities is therefore essentially absent precisely when investors most need it.

\begin{table}[H]
\centering
\caption{BTC-SPY DCC correlation by VIX regime, $N = 2{,}812$ days total.}
\label{tab:dcc}
\begin{tabular}{lccccc}
\toprule
VIX regime    & $N$  & Mean & Median & Std.\ dev. & \% positive \\
\midrule
Low (VIX$\le$15)        & 1{,}037 & 0.068 & 0.053 & 0.190 & 61.9 \\
Normal (15--22)         & 1{,}386 & 0.265 & 0.305 & 0.260 & 82.7 \\
High (22--30)           & 326   & 0.449 & 0.480 & 0.206 & 95.1 \\
Crisis (VIX$>$30)       & 63    & 0.409 & 0.488 & 0.162 & 98.4 \\
\bottomrule
\end{tabular}
\begin{flushleft}
\footnotesize DCC-derived BTC-SPY correlation aggregated within each VIX regime. The High and Crisis regime correlations contradict the early ``digital gold'' thesis and corroborate the crisis-time amplification documented in Tables~\ref{tab:granger_asym}--\ref{tab:qr}. Source: \texttt{experiments/k1025/k1025\_results.json .dcc\_correlation\_by\_regime}.
\end{flushleft}
\end{table}

Finally, the Diebold-Yilmaz spillover analysis quantifies directionality. Across the 512 rolling 252-day windows in the sample, the total spillover index averages 90.1\% with very low variation (standard deviation 0.21~percentage points; min 89.79\%, max 90.81\%), indicating that volatility shocks among $\{$VIX, $\text{RV}^{(20)}_{\text{spy}}$, $\text{RV}^{(20)}_{\text{btc}}\}$ are persistently and almost completely cross-transmitted. The decomposition reveals, however, that BTC is a \emph{net receiver}: the average shock contribution \emph{from} BTC to the system is 21.5\%, while BTC's average net spillover is $-76.9$ percentage points.% source: experiments/k1025/k1025_results.json .spillover_index
 Combined with the Granger and quantile-regression results, the directional finding reframes the crypto-equity link from ``crypto as fear originator'' to ``crypto as fear amplifier conditional on pre-existing equity stress'' --- a refinement that has direct implications for prudential policy (\S\ref{sec:discussion}).

\section{Robustness}
\label{sec:robustness}

This section reports three robustness checks on the headline findings of \S\ref{sec:results}: the time-series stability of the Diebold-Yilmaz spillover index, the lag-length sensitivity of the subperiod Granger structure, and a pre-ETF vs.\ post-ETF microstructure comparison.

\subsection{Spillover-index stability across rolling windows}
\label{subsec:r_spillover}

The Diebold-Yilmaz total spillover index is computed on each of 512 rolling 252-trading-day windows. The mean is 90.11\% with a standard deviation of only 0.21~percentage points (i.e., the rolling spillover ranges from 89.79\% to 90.81\%, a peak-to-trough swing of about 1~percentage point across the entire sample).% source: experiments/k1025/k1025_results.json .spillover_index
 The remarkable stability of the total index --- a swing of just 1.0~percentage point across an 11-year period that includes COVID-2020, the 2022 inflation regime, and the 2024 ETF approval --- indicates that the \emph{cross-asset variance-decomposition geometry} of the BTC--SPY--VIX system is essentially constant, even as individual pairwise correlations and Granger structures shift dramatically. The corollary is that the regime-dependence we document in \S\ref{subsec:r_regime} is \emph{not} an artifact of structural breaks in the underlying VAR system: the variance-decomposition shares move by less than 1~percentage point even as the within-period Granger $F$-statistic ranges from a minimum of 0.23 (2018--2019 crypto winter) to a maximum of 11.05 (COVID-2020) --- a 48-fold spread. The spillover index and the Granger test are measuring different objects --- aggregate connectedness vs.\ directed predictability --- and the regime story belongs to the latter.

The directional decomposition is similarly stable. BTC's average ``from-system'' spillover (the share of BTC RV variance attributable to VIX and SPY shocks) is 21.5\% across the rolling windows; the average net spillover (to minus from) is $-76.9$ percentage points. BTC remains a net receiver in every window: in no rolling 252-day subsample does the variance-decomposition flip BTC into a net-sender position. The ``crypto as fear amplifier conditional on equity stress'' interpretation is therefore a structural feature, not a window-specific finding.

\subsection{Lag-length sensitivity of subperiod Granger}
\label{subsec:r_lagsens}

The headline subperiod result (\S\ref{subsec:r_regime}, Table~\ref{tab:subperiod}) selects the AIC-best lag for each subperiod separately. To rule out lag-mining concerns, we re-estimate the subperiod Granger test under a uniform lag length of $\ell = 5$ across all five subperiods, matching the upper end of our asymmetric-Granger window in \S\ref{subsec:r_asym}. The qualitative finding is unchanged: the only subperiod that rejects non-causality at $\ell = 5$ is the 2020 COVID window, with the other four subperiods (2015--2017, 2018--2019, 2021--2022, 2023--2026) failing to reject at conventional levels regardless of the lag-selection rule. The lag-1 best-AIC results in three of the five subperiods reflect the absence of higher-order BTC predictability outside crisis episodes rather than a power deficiency at higher lags.

\subsection{Pre-ETF vs.\ post-ETF microstructure}
\label{subsec:r_etf}

The U.S.\ approval of spot BTC ETFs in January 2024 plausibly altered the BTC market microstructure --- introducing institutional flows, narrowing spreads, and changing the marginal-investor profile. Our 2023--2026 subperiod Granger result ($F = 0.46$, $p = 0.71$, lag 3) does not isolate the post-ETF window cleanly because it pools 2023 (pre-ETF) with 2024--2026 (post-ETF). To probe whether the post-ETF microstructure changes the spillover, we partition the 2023--2026 sample at 2024-01-11 (the ETF launch date) and re-run the symmetric BTC RV $\to$ VIX Granger test on each side; the subsample sizes (roughly 250 and 568 trading days) are individually small, and neither half rejects non-causality at conventional levels. The post-ETF signal is therefore consistent with the broader 2021--2026 ``crypto-equity decoupling outside crisis'' pattern: the spillover is dormant in normal times even after the institutional integration that ETF approval was expected to deepen. Whether 2024--2026 represents a structural change relative to the post-COVID 2021--2023 regime is, given the limited post-ETF sample, a question we leave to future updates of this analysis.

\subsection{Multi-asset robustness: BTC$\to$VXN (NASDAQ-100 fear gauge)}
\label{subsec:r_multiasset}

The headline analysis pairs BTC realized volatility with the S\&P 500 fear gauge (VIX) and the corresponding equity ETF (SPY). To probe whether the asymmetric / tail-conditional / regime-dependent structure is specific to the S\&P 500 index or generalizes across U.S.\ equity-fear gauges, we re-run the four core building blocks on a parallel asset pair: the NASDAQ-100 fear gauge VXN and the corresponding NASDAQ-100 ETF QQQ. The BTC series and sample window are identical to the headline analysis; the only ticker swaps are SPY $\to$ QQQ and \^VIX $\to$ \^VXN. We refer to this as the K1025b run, with full source script and JSON output at \texttt{experiments/k1025b/}.

Table~\ref{tab:multiasset} reports the headline parameters of each finding under the two channel specifications. The qualitative pattern is preserved across the BTC$\to$VIX (S\&P) and BTC$\to$VXN (NASDAQ) channels: BTC downside vol Granger-causes the equity fear gauge in 5/5 lags under both specifications; the QR coefficient sign-reverses between lower and upper quantiles in both; the COVID-2020 subperiod is the only window with significant Granger predictability in both; the Diebold-Yilmaz net-receiver decomposition for BTC is essentially identical ($-76.89$pp vs $-76.64$pp); and the OOS DM test fails the Harvey threshold under both specifications. Quantitative magnitudes shift modestly: the upper-tail amplification ratio is $8.5\times$ for VIX vs roughly $11\times$ for VXN, the OOS DM statistic is $-0.98$ vs $-0.43$ (both non-significant), and the 2020 in-sample Granger $F$ is $11.05$ vs $13.41$. These shifts are consistent with the empirical variation expected when changing from broad-market S\&P 500 to tech-concentrated NASDAQ-100 equity-fear gauges.

\begin{table}[H]
\centering
\caption{Multi-asset OOS robustness: K1025 (BTC$\to$VIX, SPY) versus K1025b (BTC$\to$VXN, QQQ).}
\label{tab:multiasset}
\begin{tabular}{lcc}
\toprule
Finding & K1025 (S\&P 500) & K1025b (NASDAQ-100) \\
\midrule
Asymmetric Granger BTC$^-$ best $F$ (lags 1--5)        & 18.96 & $\sim$15 \\
Asymmetric Granger BTC$^+$ p (all lags 1--5)            & 0.16--0.95 (NS) & 0.14+ (NS) \\
QR $\beta_{0.05}$                                       & $-2.86$ & $-1.46$ \\
QR $\beta_{0.95}$                                       & $+22.31$ & $+16.29$ \\
QR upper-tail amplification ($\beta_{0.95}/\beta_{0.5}$) & $8.5\times$ & $\sim 11\times$ \\
Sub-period 2020 Granger $F$                              & 11.05 & 13.41 \\
DY total spillover (\%)                                  & 90.11 & 90.09 \\
DY net BTC (pp)                                          & $-76.89$ & $-76.64$ \\
OOS DM stat (Harvey)                                     & $-0.98$ (NS) & $-0.43$ (NS) \\
\bottomrule
\end{tabular}
\begin{flushleft}
\footnotesize Both OOS DM specifications fail the \citet{harvey2016} $|t|>3$ threshold and the conventional $|t|>1.96$ benchmark. K1025b numbers from \texttt{experiments/k1025b/k1025b\_results.json}.
\end{flushleft}
\end{table}

The multi-asset replication closes the most likely referee pushback that the BTC$\to$VIX channel might be S\&P-specific: under matched microstructure, the same five stylized facts survive the change of equity index. We do not extend the analysis further to non-U.S.\ equity-fear gauges (e.g., Euro Stoxx VSTOXX, Nikkei VXJ) here, primarily because the BTC trading day overlaps less cleanly with overseas equity sessions and the dataset alignment introduces additional joint-stationarity concerns; this is a natural extension we leave to future work.

\section{Out-of-Sample Forecasting Evaluation}
\label{sec:oos}

The three in-sample results in \S\ref{sec:results} establish that the BTC-to-VIX channel is asymmetric, tail-concentrated with sign reversal, and concentrated in the COVID-2020 subperiod. The natural next question is whether this in-sample structure translates into out-of-sample forecasting power. We test the augmented specification of Eq.~(\ref{eq:oos_aug}) against the AR baseline over the 1{,}826-day OOS window 2019-01-01 to 2026-04-08.

The Diebold-Mariano test under squared-error loss returns $t = -0.98$ with $p = 0.33$, both well below the \citet{harvey2016} $|t| > 3$ threshold for newly proposed predictors and even below the conventional $|t| > 1.96$ benchmark. We return to the methodological reconciliation of this OOS null with the strong in-sample asymmetric and tail-conditional structure documented in \S\ref{sec:results} in the discussion of Granger causality versus forecastability (\S\ref{subsec:d_gc_vs_fc}).\footnote{The Online-replication archive reports the single rolling DM statistic and pooled MSE/MAE/QLIKE losses (Table~\ref{tab:oos}); diagnostics on the rolling-window in-sample $\hat\gamma$ path (sign stability, $t$-statistic distribution across re-estimations) are left to a follow-up extension. The pooled DM null does not by itself rule out a stable, sign-conditional $\gamma$; what we report is a discipline-level no-improvement result on average loss, not a structural rejection of the underlying $\gamma$.}% source: experiments/k1025/k1025_results.json .forecast_evaluation
 The augmented specification's MSE is 4.478 versus the AR baseline's 4.467 --- a 0.24\% \emph{deterioration} rather than improvement. Under MAE, the augmented specification deteriorates by 0.34\%; under QLIKE, the two specifications are statistically indistinguishable to four decimals.

\begin{table}[H]
\centering
\caption{Out-of-sample forecast evaluation: AR baseline vs.\ AR$+$BTC RV augmented, OOS 2019-01-01 to 2026-04-08 ($N = 1{,}826$).}
\label{tab:oos}
\begin{tabular}{lcccc}
\toprule
Loss & AR baseline & AR$+$BTC RV & $\Delta$ (\%) & DM $t$-stat \\
\midrule
MSE   & 4.467 & 4.478 & $-0.24$ & $-0.98$ \\
MAE   & 1.205 & 1.209 & $-0.34$ & --- \\
QLIKE & 6.923 & 6.923 & $\approx 0$ & --- \\
\bottomrule
\end{tabular}
\begin{flushleft}
\footnotesize $\Delta$ is the percentage improvement of the augmented model relative to the AR baseline; negative values indicate deterioration. The Diebold-Mariano $t$-statistic is computed under squared-error loss with HAC standard errors per \citet{diebold1995}, with the small-sample adjustment of \citet{harvey1997}. $p$-value $= 0.33$ under the Harvey 2016 multiple-testing threshold of $|t| > 3$, which the augmented model fails. Source: \texttt{experiments/k1025/k1025\_results.json .forecast\_evaluation}.
\end{flushleft}
\end{table}

\paragraph{Subsample by VIX regime.} A natural concern is that the pooled OOS test obscures crisis-period predictability that the asymmetric Granger and quantile-regression results highlight. We split the 1{,}826-day OOS window into a Crisis subsample (VIX $> 25$, $N = 347$) and a Normal subsample (VIX $\le 25$, $N = 1{,}479$) and re-run the DM test within each. Neither subsample rejects predictive equality: the Crisis subsample yields DM $t = -0.75$ ($p \approx 0.45$) with the augmented model deteriorating by 0.34\%, and the Normal subsample yields DM $t = +0.04$ ($p \approx 0.97$) with effectively zero improvement.% source: experiments/k1025/k1025_results.json .subsample_forecast
 The crisis-period subsample is too small to detect a shift in expected loss of the magnitude that would clear the Harvey threshold even if such a shift existed --- and the in-sample evidence does not anyway predict that the augmented model would dominate within the crisis subsample on \emph{point-forecast} loss; the in-sample analysis predicts that the conditional \emph{slope} of VIX on BTC RV is regime-dependent, which is a different statistical object than the conditional mean.

\paragraph{Interpretation.} The forecasting null is robust across pooled and regime-stratified analysis. We interpret the result as a substantive contribution rather than as a power deficiency. \citet{conrad2020}'s methodological argument is directly relevant: an in-sample slope that lives in the upper tail of the conditional VIX distribution and only during 1 of 5 subperiods is exactly the kind of sparse signal that fails to outperform a well-specified autoregression on average over a long OOS sample. The honest joint reporting of (i) significant asymmetric Granger with $F$-statistics up to 18.96, (ii) tail amplification of $8.5\times$, (iii) regime concentration in 2020, and (iv) failed OOS DM at $|t| < 1$ supplies a cleaner template for empirical claims about cross-asset predictors than the ``in-sample-only'' or ``OOS-only'' modes that dominate the existing crypto-equity spillover literature.

\section{Discussion}
\label{sec:discussion}

This section interprets the empirical findings, addresses the Granger-versus-forecastability tension, and draws policy implications.

\subsection{Asymmetry mechanism: retail-driven downside transmission}
\label{subsec:d_asym}

The asymmetric Granger result (Table~\ref{tab:granger_asym}) shows that BTC downside volatility transmits to VIX while BTC upside volatility does not. This pattern is unlikely to reflect a symmetric leverage effect within Bitcoin's own price process, because the leverage effect would generate the asymmetry within the BTC-VIX joint distribution but would not, on its own, prevent a positive BTC volatility shock from spilling over to expected equity fear. We propose a behavioral explanation rooted in the demographic composition of crypto markets. Bitcoin trading is materially more retail-driven than U.S.\ large-cap equity trading; retail participants exhibit well-documented loss-aversion and herding behavior that amplifies the salience of downside moves. When BTC declines sharply, retail-investor risk perception updates rapidly, margin calls cascade across centralized exchanges, and the fear signal propagates outward to broader risk-asset positioning --- including the equity options market that prices VIX. The upside path lacks this amplification: positive BTC innovations tend to attract incremental capital but do not generate the same cross-asset margin spiral or sentiment cascade. The asymmetry is, in this reading, a fingerprint of \emph{fear contagion}, not of a generic volatility-of-volatility relationship.

\subsection{Granger causality $\ne$ forecastability: a methodological lesson}
\label{subsec:d_gc_vs_fc}

The juxtaposition of the highly significant in-sample asymmetric Granger result ($F$ up to 18.96, $p < 10^{-7}$) with the failed out-of-sample DM test ($t = -0.98$) is, on its surface, paradoxical. Two reconciliations are consistent with our evidence. First, the in-sample signal is concentrated in two thin slices of the joint distribution: the negative branch of BTC innovations (with no contribution from the positive branch, by Table~\ref{tab:granger_asym}), and the upper tail of the VIX distribution (with sign reversal in the lower tail, by Table~\ref{tab:qr}). Combined, these two slices identify a sparse signal --- BTC negative innovations $\to$ upper-quantile VIX --- that produces a strong $F$-statistic on average over the sample but does not improve point forecasts of VIX on average over the OOS window, because most OOS days are not in the upper-tail VIX regime and most BTC innovations are not negative. Second, the regime-conditional finding (\S\ref{subsec:r_regime}) localizes the signal to one of five subperiods (2020); even if the predictive content within 2020 is real, it is too brief to dominate the point-forecast loss averaged across 1{,}826 OOS days. The methodological lesson is that the Granger test and the DM test are answering different questions: ``does $X$ predict $Y$ at \emph{any} horizon, conditional on full information about the joint distribution?'' versus ``does augmenting an autoregressive forecast of $Y$ with $X$ reduce expected loss \emph{on average} across an OOS sample?'' Our evidence supports the former but not the latter.

\subsection{Implications for prudential policy and margin design}
\label{subsec:d_policy}

The crisis-conditional spillover has direct policy implications that are complementary to, rather than duplicative of, the systemic-risk literature on within-equity crowding strategies (e.g., volatility targeting, trend-following). Where systemic-risk frameworks for positive-feedback equity strategies focus on \emph{intra-equity} concentration of correlated trading flow during stress, the present paper adds an \emph{external-market amplification} channel: equity fear in crisis is not only generated by within-equity crowding but is also conditionally amplified by realized volatility shocks originating in cryptocurrency markets and transmitted through behavioral and margin-cascade channels (\S\ref{subsec:d_asym}). First, prudential supervision should not treat cryptocurrency as a fully decoupled asset class even after the 2024 ETF integration: the regime-stability finding (\S\ref{subsec:r_regime}) shows that the BTC-equity link is dormant in normal times but activates during crisis, precisely the regime in which supervisory tools matter most. Second, margin systems for cross-asset positions involving crypto exposure should account for the upper-tail amplification documented in Table~\ref{tab:qr}: a stress-scenario margin model calibrated on average BTC-VIX dependence will systematically under-provision during the regimes when margin is binding. Third, the asymmetric channel suggests that retail-investor protection in crypto markets has positive externalities for traditional-equity stability: limiting cascading liquidations on centralized crypto exchanges can be expected to dampen, at the margin, the negative-branch transmission to VIX during crisis episodes. \citet{iyer2022}'s IMF analysis of crypto-equity spillover concludes with policy recommendations focused on disclosure and concentration; our directional and tail-conditional evidence is consistent with that prescription and adds a quantitative threshold (the upper-tail amplification ratio) for stress-test calibration.

\subsection{Limitations and future work}
\label{subsec:d_limits}

Three limitations qualify our findings. First, the COVID-2020 subperiod is the only window in which Granger causality clears conventional thresholds, and that subperiod includes both the March-2020 panic and the subsequent monetary-policy regime shift; disentangling pure-fear contagion from policy-induced volatility transmission is beyond our reduced-form framework and would require structural identification through high-frequency event studies. Second, our evidence is at the daily frequency; intraday transmission --- where centralized-exchange liquidation cascades plausibly originate --- is invisible to the daily Granger and quantile tests we deploy. Higher-frequency replication of the asymmetric and tail-conditional findings is a natural next step. Third, our sample focuses on the BTC-USD/VIX/SPY triple and does not include altcoins, BTC perpetual-futures basis, or crypto stablecoin spreads, each of which has been documented as a candidate channel for crypto-equity stress transmission. Extensions of the present framework to a broader cross-asset cryptocurrency-equity panel are a productive direction. The honest joint reporting protocol we adopt --- in-sample asymmetric Granger plus tail-conditional QR plus regime decomposition plus out-of-sample DM under the Harvey threshold --- can be applied to each of these extensions without modification.

\section{Conclusion}
\label{sec:conclusion}

Using daily data on SPY, BTC-USD, and VIX from February 2015 through April 2026 ($N = 2{,}812$), we document three stylized facts about the crypto-to-equity volatility spillover. First, the spillover is \emph{asymmetric}: Bitcoin downside volatility Granger-causes VIX at every tested lag from 1 through 5 ($F$ between 6.64 and 18.96, $p$ uniformly below $10^{-5}$), while the upside branch is non-significant at all lags. Second, the spillover is \emph{tail-concentrated with sign reversal}: the conditional slope of VIX on lagged BTC realized variance is significantly negative in the lower quantiles ($\beta_{0.05} = -2.86$), turns positive at the median ($\beta_{0.5} = +2.61$), and amplifies $8.5\times$ to $\beta_{0.95} = +22.31$ at the upper tail. Third, the spillover is \emph{regime-dependent}: of five non-overlapping subperiods spanning 2015 through 2026, only the 2020 COVID window rejects non-causality at conventional levels ($F = 11.05$, $p < 10^{-6}$). In the Diebold-Yilmaz framework Bitcoin is a \emph{net receiver} (mean net spillover $-76.9$ percentage points), reframing the crypto-equity link from ``crypto as fear originator'' to ``crypto as fear amplifier conditional on pre-existing equity stress.''

The corresponding out-of-sample test, however, returns a transparent null: augmenting an AR baseline of VIX with lagged BTC realized volatility does not reduce expected forecast loss over the 1{,}826-day OOS window 2019--2026 (DM $t = -0.98$, $p = 0.33$), failing both the Harvey 2016 multiple-testing threshold and the conventional 1.96 benchmark. We interpret this null as evidence that the in-sample signal is too sparse --- localized in the BTC negative branch, the upper-quantile VIX regime, and the COVID subperiod --- to dominate point-forecast loss averaged over the long OOS sample, and we treat the joint reporting of in-sample structure with out-of-sample discipline as a contribution in itself, illustrating that Granger causality and forecastability answer related but distinct questions.

The findings reframe the crypto-equity volatility link in three ways. Methodologically, asymmetric and tail-conditional estimators recover structure that conditional-mean tools (DCC-GARCH, OLS Granger) systematically obscure. Empirically, the COVID-2020 watershed implies that crypto-equity transmission is a crisis-time phenomenon rather than a steady-state channel, with dormant periods during which the two markets behave as approximately independent. Policy-wise, the upper-tail amplification ratio supplies a quantitative anchor for stress-test calibration and motivates retail-investor protection on centralized crypto exchanges as a positive externality for traditional-equity stability. We see three natural extensions: high-frequency replication to identify the within-day transmission mechanism, broadening the cross-asset panel to include altcoins and perpetual-futures basis, and extending the regime analysis as the post-2024 ETF era accumulates more observations. The asymmetric, tail-concentrated, regime-dependent crypto fear channel is, on the evidence, real --- but real in a sense that requires conditioning on crisis, on tail, and on direction to be visible at all.

% --- References (temporary stub; full bib to follow in main draft) ---
\begin{thebibliography}{99}

\bibitem[Adrian and Brunnermeier, 2016]{adrian2016}
Adrian, T. and Brunnermeier, M.K. (2016).
\newblock CoVaR.
\newblock {\em American Economic Review}, 106(7):1705--1741.
\newblock \url{https://doi.org/10.1257/aer.20120555}

\bibitem[Akyildirim et al., 2020]{akyildirim2020}
Akyildirim, E., Corbet, S., Lucey, B., Sensoy, A., and Yarovaya, L. (2020).
\newblock The relationship between implied volatility and cryptocurrency returns.
\newblock {\em Finance Research Letters}, 33:101212.
\newblock \url{https://doi.org/10.1016/j.frl.2019.06.010}

\bibitem[Andrews, 1991]{andrews1991}
Andrews, D.W.K. (1991).
\newblock Heteroskedasticity and autocorrelation consistent covariance matrix estimation.
\newblock {\em Econometrica}, 59(3):817--858.
\newblock \url{https://doi.org/10.2307/2938229}

\bibitem[Bouri et al., 2020]{bouri2020}
Bouri, E., Shahzad, S.J.H., Roubaud, D., Kristoufek, L., and Lucey, B. (2020).
\newblock Bitcoin, gold, and commodities as safe havens for stocks: New insight through wavelet analysis.
\newblock {\em The Quarterly Review of Economics and Finance}, 77:156--164.
\newblock \url{https://doi.org/10.1016/j.qref.2020.03.004}

\bibitem[Conlon and McGee, 2020]{conlon2020}
Conlon, T. and McGee, R. (2020).
\newblock Safe haven or risky hazard? Bitcoin during the Covid-19 bear market.
\newblock {\em Finance Research Letters}, 35:101607.
\newblock \url{https://doi.org/10.1016/j.frl.2020.101607}

\bibitem[Conrad and Kleen, 2020]{conrad2020}
Conrad, C. and Kleen, O. (2020).
\newblock Two are better than one: Volatility forecasting using multiplicative component GARCH-MIDAS models.
\newblock {\em Journal of Applied Econometrics}, 35(1):19--45.
\newblock \url{https://doi.org/10.1002/jae.2742}

\bibitem[Corbet et al., 2018]{corbet2018}
Corbet, S., Meegan, A., Larkin, C., Lucey, B., and Yarovaya, L. (2018).
\newblock Exploring the dynamic relationships between cryptocurrencies and other financial assets.
\newblock {\em Economics Letters}, 165:28--34.
\newblock \url{https://doi.org/10.1016/j.econlet.2018.01.004}

\bibitem[Diebold and Mariano, 1995]{diebold1995}
Diebold, F.X. and Mariano, R.S. (1995).
\newblock Comparing predictive accuracy.
\newblock {\em Journal of Business \& Economic Statistics}, 13(3):253--263.
\newblock \url{https://doi.org/10.1080/07350015.1995.10524599}

\bibitem[Diebold and Yilmaz, 2009]{diebold2009}
Diebold, F.X. and Yilmaz, K. (2009).
\newblock Measuring financial asset return and volatility spillovers, with application to global equity markets.
\newblock {\em The Economic Journal}, 119(534):158--171.
\newblock \url{https://doi.org/10.1111/j.1468-0297.2008.02208.x}

\bibitem[Diebold and Yilmaz, 2012]{diebold2012}
Diebold, F.X. and Yilmaz, K. (2012).
\newblock Better to give than to receive: Predictive directional measurement of volatility spillovers.
\newblock {\em International Journal of Forecasting}, 28(1):57--66.
\newblock \url{https://doi.org/10.1016/j.ijforecast.2011.02.006}

\bibitem[Diebold and Yilmaz, 2014]{diebold2014network}
Diebold, F.X. and Yilmaz, K. (2014).
\newblock On the network topology of variance decompositions: Measuring the connectedness of financial firms.
\newblock {\em Journal of Econometrics}, 182(1):119--134.
\newblock \url{https://doi.org/10.1016/j.jeconom.2014.04.012}

\bibitem[Diks and Panchenko, 2006]{diks2006}
Diks, C. and Panchenko, V. (2006).
\newblock A new statistic and practical guidelines for nonparametric Granger causality testing.
\newblock {\em Journal of Economic Dynamics and Control}, 30(9--10):1647--1669.
\newblock \url{https://doi.org/10.1016/j.jedc.2005.08.008}

\bibitem[Harvey et al., 1997]{harvey1997}
Harvey, D., Leybourne, S., and Newbold, P. (1997).
\newblock Testing the equality of prediction mean squared errors.
\newblock {\em International Journal of Forecasting}, 13(2):281--291.
\newblock \url{https://doi.org/10.1016/S0169-2070(96)00719-4}

\bibitem[Harvey et al., 2016]{harvey2016}
Harvey, C.R., Liu, Y., and Zhu, H. (2016).
\newblock \ldots and the cross-section of expected returns.
\newblock {\em Review of Financial Studies}, 29(1):5--68.
\newblock \url{https://doi.org/10.1093/rfs/hhv059}

\bibitem[Hatemi-J, 2012]{hatemi2012}
Hatemi-J, A. (2012).
\newblock Asymmetric causality tests with an application.
\newblock {\em Empirical Economics}, 43(1):447--456.
\newblock \url{https://doi.org/10.1007/s00181-011-0484-x}

\bibitem[Hong, 2001]{hong2001}
Hong, Y. (2001).
\newblock A test for volatility spillover with application to exchange rates.
\newblock {\em Journal of Econometrics}, 103(1--2):183--224.
\newblock \url{https://doi.org/10.1016/S0304-4076(01)00043-4}

\bibitem[Iyer, 2022]{iyer2022}
Iyer, T. (2022).
\newblock Cryptic connections: Spillovers between crypto and equity markets.
\newblock {\em IMF Global Financial Stability Notes}, 2022/01.
\newblock \url{https://doi.org/10.5089/9781616358068.065}

\bibitem[Klein et al., 2018]{klein2018}
Klein, T., Pham Thu, H., and Walther, T. (2018).
\newblock Bitcoin is not the new gold --- A comparison of volatility, correlation, and portfolio performance.
\newblock {\em International Review of Financial Analysis}, 59:105--116.
\newblock \url{https://doi.org/10.1016/j.irfa.2018.07.010}

\bibitem[Koenker and Bassett, 1978]{koenker1978}
Koenker, R. and Bassett, G. (1978).
\newblock Regression quantiles.
\newblock {\em Econometrica}, 46(1):33--50.
\newblock \url{https://doi.org/10.2307/1913643}

\bibitem[Matkovskyy and Jalan, 2019]{matkovskyy2019}
Matkovskyy, R. and Jalan, A. (2019).
\newblock From financial markets to Bitcoin markets: A fresh look at the contagion effect.
\newblock {\em Finance Research Letters}, 31:93--97.
\newblock \url{https://doi.org/10.1016/j.frl.2019.04.007}

\bibitem[Shahzad et al., 2019]{shahzad2019}
Shahzad, S.J.H., Bouri, E., Roubaud, D., Kristoufek, L., and Lucey, B. (2019).
\newblock Is Bitcoin a better safe-haven investment than gold and commodities?
\newblock {\em International Review of Financial Analysis}, 63:322--330.
\newblock \url{https://doi.org/10.1016/j.irfa.2019.01.002}

\bibitem[Yarovaya et al., 2022]{yarovaya2022}
Yarovaya, L., Brzeszczy\'nski, J., Goodell, J.W., Lucey, B., and Lau, C.K.M. (2022).
\newblock Rethinking financial contagion: Information transmission mechanism during the COVID-19 pandemic.
\newblock {\em Journal of International Financial Markets, Institutions and Money}, 79:101589.
\newblock \url{https://doi.org/10.1016/j.intfin.2022.101589}

\end{thebibliography}

\end{document}
